% GENERATED FILE. Edit canonical lessons, then run npm run generate:books.
% canonical-source-hash: fnv1a64:1188ba72fdd310d3
% canonical-lessons: GE-C220-puenktlich, GE-C220-regnerisch, GE-C220-roh, GE-C220-salzig, GE-C220-sauer

\chapter{Punctual, Rainy, Raw, Salty, Sour}
\label{ch:ge-punctual-rainy-raw-salty-sour}
\hlchaptermodality{220}

\begin{chapteropening}
\textbf{By the end of this chapter:} \emph{I can say punctual, rainy, raw, salty and sour.}

Complete the last lesson of chapter 220: I can say punctual, rainy, raw, salty and sour.
\end{chapteropening}

\section[pünktlich]{pünktlich — punctual}

\label{lesson:GE-C220-puenktlich}



\begin{quote}
Before the new one: say the German for tidy, proper, then the German for personal.
\end{quote}

\subsection*{You'll want to know: pünktlich}
\textbf{pünktlich} — \textquotedblleft{}punctual\textquotedblright{}.

\begin{cousinweb}
\textbf{pünktlich} is the word for \textquotedblleft{}point\textquotedblright{}, from Latin \emph{pūnctum}, \textquotedblleft{}a prick, a point\textquotedblright{} (EtymDB 2.1), + \emph{-lich}: on the dot.
\end{cousinweb}

\begin{grammarlens}[title={What you've built}]
One more word for this chapter.
\end{grammarlens}

\subsection*{Your turn}
\begin{itemize}
\raggedright
  \item \emph{Say it:} \emph{pünktlich}
  \item \emph{Say it:} \emph{pünktlich}, once more
  \item \emph{Say it:} say \emph{pünktlich}
  \item \emph{Recall:} say the German for tidy, proper, then the German for personal, then say \emph{pünktlich} again
\end{itemize}

\begin{tcolorbox}[breakable,colback=teal!4,colframe=teal!35!black,title={Before you move on}]
What does pünktlich mean? (\textquotedblleft{}Punctual\textquotedblright{}.) Say it once more.
\end{tcolorbox}
\section[regnerisch]{regnerisch — rainy}

\label{lesson:GE-C220-regnerisch}



\begin{quote}
Before the new one: say the German for personal, then the German for punctual.
\end{quote}

\subsection*{You'll want to know: regnerisch}
\textbf{regnerisch} — \textquotedblleft{}rainy\textquotedblright{}.

\begin{cousinweb}
\textbf{regnerisch} is built on \emph{Regen}, \textquotedblleft{}rain\textquotedblright{}, + \emph{-isch}. \emph{Regen} $\leftarrow$ Proto-Germanic \emph{*regnaz}, English \emph{rain}. (EtymDB 2.1)
\end{cousinweb}

\begin{grammarlens}[title={What you've built}]
One more word for this chapter.
\end{grammarlens}

\subsection*{Your turn}
\begin{itemize}
\raggedright
  \item \emph{Say it:} \emph{regnerisch}
  \item \emph{Say it:} \emph{regnerisch}, once more
  \item \emph{Say it:} say \emph{regnerisch}
  \item \emph{Recall:} say the German for personal, then the German for punctual, then say \emph{regnerisch} again
\end{itemize}

\begin{tcolorbox}[breakable,colback=teal!4,colframe=teal!35!black,title={Before you move on}]
What does regnerisch mean? (\textquotedblleft{}Rainy\textquotedblright{}.) Say it once more.
\end{tcolorbox}
\section[roh]{roh — raw}

\label{lesson:GE-C220-roh}



\begin{quote}
Before the new one: say the German for punctual, then the German for rainy.
\end{quote}

\subsection*{You'll want to know: roh}
\textbf{roh} — \textquotedblleft{}raw\textquotedblright{}.

\begin{cousinweb}
\textbf{roh} $\leftarrow$ Middle High German \emph{rō} $\leftarrow$ Old High German \emph{rāo} — inherited, the same word as English \emph{raw}. (EtymDB 2.1)
\end{cousinweb}

\begin{grammarlens}[title={What you've built}]
One more word for this chapter.
\end{grammarlens}

\subsection*{Your turn}
\begin{itemize}
\raggedright
  \item \emph{Say it:} \emph{roh}
  \item \emph{Say it:} \emph{roh}, once more
  \item \emph{Say it:} say \emph{roh}
  \item \emph{Recall:} say the German for punctual, then the German for rainy, then say \emph{roh} again
\end{itemize}

\begin{tcolorbox}[breakable,colback=teal!4,colframe=teal!35!black,title={Before you move on}]
What does roh mean? (\textquotedblleft{}Raw\textquotedblright{}.) Say it once more.
\end{tcolorbox}
\section[salzig]{salzig — salty}

\label{lesson:GE-C220-salzig}



\begin{quote}
Before the new one: say the German for rainy, then the German for raw.
\end{quote}

\subsection*{You'll want to know: salzig}
\textbf{salzig} — \textquotedblleft{}salty\textquotedblright{}.

\begin{cousinweb}
\textbf{salzig} is \emph{Salz}, \textquotedblleft{}salt\textquotedblright{}, + \emph{-ig}, built like English \emph{salty}. \emph{Salz} $\leftarrow$ Proto-Germanic \emph{*saltą}. (EtymDB 2.1)
\end{cousinweb}

\begin{grammarlens}[title={What you've built}]
One more word for this chapter.
\end{grammarlens}

\subsection*{Your turn}
\begin{itemize}
\raggedright
  \item \emph{Say it:} \emph{salzig}
  \item \emph{Say it:} \emph{salzig}, once more
  \item \emph{Say it:} say \emph{salzig}
  \item \emph{Recall:} say the German for rainy, then the German for raw, then say \emph{salzig} again
\end{itemize}

\begin{tcolorbox}[breakable,colback=teal!4,colframe=teal!35!black,title={Before you move on}]
What does salzig mean? (\textquotedblleft{}Salty\textquotedblright{}.) Say it once more.
\end{tcolorbox}
\section[sauer]{sauer — sour, annoyed}

\label{lesson:GE-C220-sauer}



\begin{quote}
Before the new one: say the German for raw, then the German for salty.
\end{quote}

\subsection*{You'll want to know: sauer}
\textbf{sauer} — \textquotedblleft{}sour, annoyed\textquotedblright{}.

\begin{cousinweb}
\textbf{sauer} $\leftarrow$ Middle High German \emph{sūr} $\leftarrow$ Old High German \emph{sūr} (EtymDB 2.1) — inherited, the same word as English \emph{sour}. \textquotedblleft{}Annoyed\textquotedblright{} is a soured mood.
\end{cousinweb}

\begin{grammarlens}[title={What you've built}]
One more word for this chapter.
\end{grammarlens}

\subsection*{Your turn}
\begin{itemize}
\raggedright
  \item \emph{Say it:} \emph{sauer}
  \item \emph{Say it:} \emph{sauer}, once more
  \item \emph{Say it:} say \emph{sauer}
  \item \emph{Recall:} say the German for raw, then the German for salty, then say \emph{sauer} again
\end{itemize}

\begin{tcolorbox}[breakable,colback=teal!4,colframe=teal!35!black,title={Before you move on}]
What does sauer mean? (\textquotedblleft{}Sour, annoyed\textquotedblright{}.) Say it once more.
\end{tcolorbox}
